\scanpage{7}
image inverse de l'extension $(5_I)$, par l'homomorphisme (4) de $\underbrace{\{\pm1\}}_{\mu_2}$ dans $U^I$.

Je vais donner une description bcp plus simple de cette extension, en \deleted{\unsure{notant que}}\inserted{factorisant} comme en (4)
\[
\mu_2\xrightarrow{\mathrm{diag}}\mu_2^I\longrightarrow U^I.
\]
On est donc réduit à déterminer l'ext. de $\mu_2$ par $\mathbb Z$ déduite de l'ext. can. (5) de $U$ par $\mathbb Z$, par le plongement
\[
\mu_2\subset U
\]
extension qui est canoniquement isom. à l'ext. \unsure{standart}

$(14)$\quad
\diagram{\begin{tikzcd}[ampersand replacement=\&, row sep=small]
0\arrow[r,no head]\&\mathbb Z\arrow[r,"2"]\&\mathbb Z\arrow[r]\&\mathbb Z/2\mathbb Z\arrow[r]\arrow[d,phantom,"\wr"]\&0\\
\&\&\&\mu_2\&
\end{tikzcd}}

Donc on prend la puissance cartésienne $I$. ième de cette extension

\diagram{\begin{tikzcd}[ampersand replacement=\&, row sep=small]
0\arrow[r,no head]\&\mathbb Z^I\arrow[r,"2"]\&\mathbb Z^I\arrow[r]\&(\mathbb Z/2\mathbb Z)^I\arrow[r,no head]\arrow[d,phantom,"\wr"]\&0\\
\&\&\&\mu_2^I\&
\end{tikzcd}}

qui donne l'image inverse de $(5_I)$ par $\mu_2^I\;-\;U^I$, puis on prend l'ext induite (de $\mu_2$ par $\mathbb Z^I$) déduite de l'hom diagonal $\mu_2\;-\;\mu_2^I$. On peut considérer l'hom d'extensions

\diagram{\begin{tikzcd}[ampersand replacement=\&, column sep=small,row sep=large]
0\arrow[r]\&\mathbb Z\arrow[r,"2"]\arrow[d,"\mathrm{diag}"']\&\mathbb Z\arrow[r,no head]\arrow[d,"\mathrm{diag}"']\&\mathbb Z/2\mathbb Z\arrow[r]\arrow[d,"\mathrm{diag}"]\&0\\
0\arrow[r]\&\mathbb Z^I\arrow[r,"2"]\&\mathbb Z^I\arrow[r]\&(\mathbb Z/2\mathbb Z)^I\arrow[r,no head]\&0
\end{tikzcd}}
\mspage{6}
